\bigskip
\noindent Autres pistes de l'onglet Findings pour ce thème, non démontrées ici :
\begin{itemize}\small\setlength{\itemsep}{0pt}
\item Cote~8 : quand $\pi A$ a des puissances divisées, le relèvement de Lubin--Tate existe et est unique par la seule théorie des déformations cristallines ; sans puissances divisées, l'argument donne l'unicité à isogénie près. La lecture de la page est à vérifier sur le fac-similé. Mathlib n'a pas la théorie des déformations cristallines. On se reportera aux notes de Grothendieck (\lot{8}{3}{43}).\piste{8-lubin-tate-via-hodge-filtration}
\item Cote~8 : deux groupes de Barsotti--Tate de dimension un et de hauteur $[K : \mathbf Q_p]$ sur $\mathbf F_q$, munis d'une action de $A$, sont isogènes si et seulement si $\varpi = \varpi'$, et $\mathrm{Hom}_A(G_0, G'_0)$ est alors explicite. La lecture de la page est à vérifier sur le fac-similé. Mathlib n'a pas les groupes de Barsotti--Tate. On se reportera aux notes de Grothendieck (\lot{8}{2}{39}).\piste{8-dim-one-A-module-homs}
\end{itemize}

\begin{encadre}
\begin{theoreme*}[Grothendieck, pages 1 et 4 à 7, valeurs propres unités de la cohomologie évanescente\piste{4-vanishing-unit-root-bound}]\nonverifie
On admet le théorème de Lefschetz faible cristallin avec torsion. Pour une section hyperplane lisse $Y$ d'une variété projective lisse $X$ de dimension $n$ sur un corps parfait de caractéristique $p$, le nombre de valeurs propres du Frobenius qui sont des unités $p$-adiques sur la cohomologie cristalline évanescente est au moins le rang $F$-semi-simple de la cohomologie cohérente évanescente $\operatorname{Coker}(H^{n-1}(X, \mathcal O_X) \to H^{n-1}(Y, \mathcal O_Y))$, le défaut étant mesuré par la $p$-torsion de $H^n_{\mathrm{cris}}(X)$.
\end{theoreme*}
\end{encadre}
{\small Les feuillets en esquissent une démonstration incomplète (pages 1 et 4 à 6). Le titre dit l'estimation \q{heuristique}, et l'énoncé~8.4.4~a) de la page~1 y est sans argument. Mathlib n'a pas la cohomologie cristalline. On se reportera aux notes de Grothendieck (\lot{4}{1}{1}).\par}

\begin{encadre}
\begin{theoreme*}[Grothendieck, pages 52 à 56, modules des variétés abéliennes ordinaires\piste{6-ordinary-moduli-complete-intersection}]\nonverifie
Pour une polarisation de degré quelconque, divisible par $p$ compris, l'anneau local en un point du schéma de modules des variétés abéliennes ordinaires polarisées en caractéristique~$p$ (degré de polarisation fixé, structure de niveau de Jacobi) est d'intersection complète et localement irréductible, et le schéma réduit est lisse : les modules formels des couples d'extensions compatibles forment le noyau d'un homomorphisme $\alpha - \beta$ de tores formels.
\end{theoreme*}
\end{encadre}
{\small Les feuillets en esquissent une démonstration incomplète : ils donnent le noyau de $\alpha - \beta$, mais le texte des pages 52 à 56 est en grande partie illisible et reconstitué, et l'hypothèse que $M_0$ est étale est ajoutée par l'édition.\par}

\begin{encadre}
\begin{theoreme*}[Grothendieck, pages 49 et 50, un groupe formel hors des variétés abéliennes\piste{6-formal-group-not-in-abelian-variety}]\nonverifie
Soient $k$ un corps fini et $E$ une courbe elliptique supersingulière sur $k$ dont tous les endomorphismes sont définis sur~$k$. La torsion du groupe formel de $E$ par une unité $\theta$ de $\mathrm{End}(E) \otimes \mathbf Z_p$ de norme réduite transcendante donne une forme de $G_{1,1}$ qui ne se plonge dans le groupe formel d'aucune variété abélienne sur $k$ ni sur aucune extension finie de~$k$.
\end{theoreme*}
\end{encadre}
{\small Les feuillets en esquissent une démonstration incomplète : ils donnent l'obstruction, d'après Weil, et disent \q{facile} l'existence de~$\theta$, qu'ils ne montrent pas.\par}

\begin{encadre}
\begin{theoreme*}[Grothendieck, pages 47 et 48, même groupe formel, non isogènes\piste{6-same-formal-group-not-isogenous}]\nonverifie
Sur un corps algébriquement clos de caractéristique~$p$, il existe des surfaces abéliennes simples de même groupe formel $\hat{\mathbf G}_m \times G_{1,1}$ qui ne sont pas isogènes.
\end{theoreme*}
\end{encadre}
{\small Les feuillets en esquissent une démonstration incomplète : la discrétude de la classe d'isogénie n'y est pas justifiée, et la page compte une famille \q{de dim 2 au plus}, chiffre de lecture douteuse, là où il faut~$1$.\par}

\begin{encadre}
\begin{theoreme*}[Grothendieck, pages 39 à 44 et 48 à 56, propriété universelle des complexes libres\piste{137-perfect-complexes-universal-sigma-exact}]\nonverifie
Soient $\mathcal P$ la catégorie $k$-linéaire des complexes bornés de $k$-modules libres de type fini, $\Sigma$ les suites exactes courtes scindées degré par degré et $\mathbb C^\bullet$ le complexe des disques. Pour toute catégorie $\mathcal B$ $k$-additive karoubienne munie d'une classe $\Sigma$ de suites exactes où tout morphisme a un $\Sigma$-noyau, $F \mapsto F(\mathbb C^\bullet)$ est une équivalence entre les foncteurs $k$-linéaires $\mathcal P \to \mathcal B$ qui envoient $\Sigma$ sur des suites exactes à gauche et les complexes de cochaînes de $\mathcal B$ ; les foncteurs $\Sigma$-exacts y correspondent aux $\Sigma$-résolutions.
\end{theoreme*}
\end{encadre}
{\small L'encadré suit la forme de la page~55 ; la page~50 hésite et biffe \q{équivalence} pour \q{pleinement fidèle}. Les feuillets en esquissent une démonstration incomplète : les propositions 1 à 3 et la preuve des pages 52 à 55 sont en formules, et la page~52 note en marge \q{il faut condition de platitude ?}. La pleine fidélité, première moitié de a) et forme retenue par la page~50, est vérifiée en Lean et énoncée ci-dessous ; l'essentielle surjectivité (seconde moitié de a), là où servent les noyaux de $\mathcal B$) l'est aussi, pour $\mathcal B$ munie de noyaux et de sommes directes finies, dans le troisième encadré ; la forme karoubienne à classe $\Sigma$ de l'énoncé et b) ne le sont pas.\par}

\begin{encadre}
\begin{theoreme*}[Grothendieck, pages 40 et 50 à 55, la pleine fidélité\piste{137-perfect-complexes-universal-sigma-exact}]\verifie{Folder137}
Soient $k$ un anneau commutatif, $\mathcal P$ la catégorie des complexes de chaînes de $k$-modules libres de type fini, en degrés $\geqslant 0$ et nuls en grand degré, $\Sigma$ la classe des suites $0 \to L' \to L \to L'' \to 0$ exactes en chaque degré, et $\mathbb C^\bullet$ le complexe des disques, $\mathbb C^n = (k \xrightarrow{\mathrm{id}} k)$ en degrés $n+1$, $n$. Alors :
\begin{enumerate}
\item (proposition~1, page~40) $\mathrm{End}(\mathbb C^n) = k \cdot \mathrm{id}$, $\mathrm{Hom}(\mathbb C^n, \mathbb C^{n+1}) = k \cdot d^n$, et $\mathrm{Hom}(\mathbb C^n, \mathbb C^j) = 0$ pour $j \notin \{n, n+1\}$ ; en particulier $d^{n+1} d^n = 0$ ;
\item pour toute catégorie $k$-linéaire $\mathcal B$, sans hypothèse de noyaux, le foncteur $F \mapsto F(\mathbb C^\bullet)$, des foncteurs $k$-linéaires $\mathcal P \to \mathcal B$ qui envoient chaque suite de $\Sigma$ sur un diagramme noyau $0 \to F(L') \to F(L) \to F(L'')$ vers les complexes de cochaînes de $\mathcal B$, est pleinement fidèle.
\end{enumerate}
\end{theoreme*}
\end{encadre}
{\small \emph{Démonstration.} (1) Un morphisme $f : \mathbb C^n \to X$ est déterminé par l'image du générateur $1$ en degré $n+1$, sa composante de degré $n$ s'en déduisant par la différentielle ; dans $\mathbb C^j$ cette image vit dans $(\mathbb C^j)_{n+1}$, qui vaut $k$ pour $j = n$ (sommet) ou $j = n+1$ (base) et $0$ sinon. (2) Tout $L$ s'insère dans une suite de $\Sigma$
\[
0 \to L \xrightarrow{\iota} I(L) \xrightarrow{\pi} L[1] \to 0, \qquad I(L)_m = L_m \oplus L_{m-1},\quad d(a, b) = (b, 0),\quad \iota(a) = (a, da),\quad \pi(a, b) = b - da,
\]
où $I(L)$, le cône de l'identité, est somme finie de disques : si $(b)$ parcourt une base de chaque $L_n$, $\sum_{n, b} p_{n,b}\, e_{n,b} = \mathrm{id}_{I(L)}$, avec $e_{n,b} : \mathbb C^n \to I(L)$ envoyant le générateur sur $(0, b)$ et $p_{n,b} : I(L) \to \mathbb C^n$ donné par la coordonnée $b^*$. Soient $F$, $G$ deux foncteurs comme dans l'énoncé et $\varphi : F(\mathbb C^\bullet) \to G(\mathbb C^\bullet)$ un morphisme de complexes. Par (1), $\varphi$ commute à toute flèche entre disques, $c\cdot\mathrm{id}$, $c \cdot d^n$ ou $0$, puisque $F$ et $G$ sont $k$-linéaires. On pose $\Phi_{I(L)} = \sum_{n,b} G(e_{n,b})\,\varphi_n\,F(p_{n,b})$ ; en insérant les deux décompositions de l'identité, $\Phi$ commute à toute flèche $I(L) \to I(M)$ et à toute flèche $\mathbb C^n \to I(M)$. Comme $G$ envoie la suite ci-dessus sur un diagramme noyau, $G(\iota)$ est un monomorphisme noyau de $G(\pi)$ ; la flèche $\Phi_{I(L)} F(\iota)$ est tuée par $G(\pi)$, car après le monomorphisme $G(\iota_{L[1]})$ elle devient $\Phi_{I(L[1])} F(\iota_{L[1]} \pi \iota) = 0$. Elle se relève donc en $\eta_L : F(L) \to G(L)$. La naturalité de $\eta$ vient de $\iota_M g = I(g)\,\iota_L$, où $I(g) = g \oplus g$ ; on a $\eta_{\mathbb C^n} = \varphi_n$ en appliquant la commutation de $\Phi$ à $\iota : \mathbb C^n \to I(\mathbb C^n)$. C'est la plénitude. La fidélité suit de même : une transformation $\eta$ vaut $\Phi$ sur $I(L)$ par la décomposition, puis est déterminée sur $L$ par le monomorphisme $G(\iota)$. Le doute marginal de la page~52 (\q{il faut condition de platitude ?}) ne se pose pas pour cette moitié : on n'y utilise ni noyaux dans $\mathcal B$, ni platitude, ni la caractérisation par récurrence sur la longueur des pages 53 et 54, que la suite du cône remplace. Lean (\href{\leanblob/Folder137.lean}{\texttt{Folder137.\allowbreak{}lean}}) définit $\mathcal P$ (\texttt{Perf}), les disques (\texttt{disc}, \texttt{dmap}, \texttt{univ}), $\Sigma$ (\texttt{IsSigma}) et les foncteurs $\Sigma$-exacts à gauche (\texttt{SigmaLeftExact}, \texttt{Lex}), et démontre (1) (\texttt{disc\_endo}, \texttt{disc\_succ}, \texttt{disc\_zero}), la suite du cône et sa décomposition (\texttt{isSigma\_cone}, \texttt{total}) et (2) (\texttt{\(\rho\)FullyFaithful}, avec \texttt{discNat}, \texttt{phi\_nat}, \texttt{extend}). L'essentielle surjectivité fait l'objet de l'encadré suivant ; ne sont pas démontrés la partie b), les propositions 2 et 3 (pages 42 à 44) et la version cosimpliciale des pages 45 et 46.\par}

\begin{encadre}
\begin{theoreme*}[Grothendieck, pages 50 à 55, l'essentielle surjectivité\piste{137-perfect-complexes-universal-sigma-exact}]\verifie{Folder137}
Avec les notations de l'encadré précédent, soit $\mathcal B$ une catégorie $k$-linéaire munie de noyaux et de sommes directes finies. Pour tout complexe de cochaînes $C^\bullet$ de $\mathcal B$, il existe un foncteur $k$-linéaire $\sigma(C^\bullet) : \mathcal P \to \mathcal B$, qui envoie chaque suite de $\Sigma$ sur un diagramme noyau, tel que $\sigma(C^\bullet)(\mathbb C^\bullet) \simeq C^\bullet$. Par conséquent $F \mapsto F(\mathbb C^\bullet)$ est une équivalence de catégories.
\end{theoreme*}
\end{encadre}
{\small \emph{Démonstration.} Pour un objet $T$ de $\mathcal B$, $M^\bullet = \mathrm{Hom}(T, C^\bullet)$ est un complexe de $k$-modules ; soit $\Phi_T(L) \subset \prod_n L_n \otimes_k M^n$ l'ensemble des familles $(z_n)$ telles que $(1 \otimes d)\, z_n = (d \otimes 1)\, z_{n+1}$ dans $L_n \otimes M^{n+1}$. Comme $L_n \otimes M = \mathrm{Hom}_k(L_n^\vee, M)$ pour $L_n$ libre de type fini, ce sont les morphismes de complexes $L^\vee \to M^\bullet$, c'est-à-dire les cycles de degré $0$ de $\underline{\mathrm{Hom}}^\bullet_k(L^\bullet, C^\bullet)$ vus à travers $T$ ; on définit $\sigma(C^\bullet)(L)$ comme un objet qui représente $T \mapsto \Phi_T(L)$. (i) $\Phi_T$ est exact à gauche sur $\Sigma$ : une suite de $\Sigma$ est scindée en chaque degré, $L''_n$ étant libre, donc reste exacte après $- \otimes_k M^n$ ; un élément de $\Phi_T(L)$ tué dans $\Phi_T(L'')$ provient d'un unique élément de $\Phi_T(L')$, la condition de cycle se transportant par injectivité. (ii) Le disque $\mathbb C^n$ est représenté par $C^n$ : un cycle de $\mathbb C^n$ est déterminé par sa composante de degré $n$, $x \in \mathrm{Hom}(T, C^n)$, celle de degré $n+1$ valant $dx$ ; par cette identification, $d^n$ correspond à la différentielle de $C^\bullet$. (iii) $I(L)$ est représenté par la somme directe finie $\bigoplus_{n,b} C^n$, car à la relation $\sum_{n,b} p_{n,b}\, e_{n,b} = \mathrm{id}$ s'ajoute l'orthogonalité $p_{m,c}\, e_{n,b} = \delta_{(n,b),(m,c)}$. (iv) $L$ est représenté par le noyau de la flèche $\sigma(I(L)) \to \sigma(I(L[1]))$ qui représente $I(L) \xrightarrow{\pi} L[1] \xrightarrow{\iota} I(L[1])$, par (i) appliqué aux suites du cône de $L$ et de $L[1]$. Les flèches de $\sigma(C^\bullet)$ s'obtiennent par Yoneda ; la $k$-linéarité et l'exactitude à gauche sur $\Sigma$ se lisent sur $\Phi$, et $\sigma(C^\bullet)(\mathbb C^n) \simeq C^n$, compatiblement aux différentielles, parce que deux objets qui représentent le même foncteur sont isomorphes. Il n'y a pas de récurrence sur la longueur ni de platitude. Les sommes directes finies sont nécessaires en plus des noyaux que nomme la page~50 : un foncteur $k$-linéaire envoie $\mathbb C^0 \oplus \mathbb C^0$ sur une somme directe $F(\mathbb C^0) \oplus F(\mathbb C^0)$, qui n'existe pas dans la sous-catégorie pleine $\{0, k\}$ des $k$-espaces vectoriels ($k$ un corps), laquelle a des noyaux ; le complexe réduit à $k$ en degré $0$ n'y est donc pas de la forme $F(\mathbb C^\bullet)$ (argument de l'édition, non formalisé). Lean (\href{\leanblob/Folder137.lean}{\texttt{Folder137.\allowbreak{}lean}}) définit $\Phi$ (\texttt{Phi}, \texttt{PhiMap}, \texttt{PhiPre}) et la représentabilité (\texttt{Rep}), et démontre (i) (\texttt{PhiMap\_injective}, \texttt{PhiMap\_exact}), (ii) (\texttt{discRep}, \texttt{discRep\_d}), (iii) (\texttt{eDec\_pDec\_self}, \texttt{eDec\_pDec\_ne}, \texttt{coneRep}) et (iv) (\texttt{kerRep}) ; il en tire $\sigma(C^\bullet)$ (\texttt{\(\sigma\)Fun}, \texttt{\(\sigma\)\_sigmaLeftExact}, \texttt{\(\sigma\)Lex}), l'isomorphisme $\sigma(C^\bullet)(\mathbb C^\bullet) \simeq C^\bullet$ (\texttt{\(\sigma\)Iso}) et, avec \texttt{\(\rho\)FullyFaithful}, l'équivalence (\texttt{\(\rho\)Equivalence}). Ne sont pas démontrés : la forme karoubienne à classe $\Sigma$ de suites exactes de $\mathcal B$ du premier encadré, la partie b), les propositions 2 et 3 (pages 42 à 44) et la version cosimpliciale des pages 45 et 46.\par}

\begin{encadre}
\begin{conjecture*}[Grothendieck, pages 16 à 18, un site $p$-adique sans puissances divisées\piste{7-noninfinitesimal-site}]
Soit $X_0$ un schéma sur un corps parfait $k$ de caractéristique $p$. La cohomologie du site des $W$-schémas formels affines $p$-adiquement complets $\operatorname{Spf} B$ munis d'une surjection $B \to A_0$ sur l'anneau d'un ouvert affine de $X_0$, de noyau nilpotent modulo $p$, sans puissances divisées ni surconvergence, est une cohomologie $p$-adique : elle redonne la cohomologie de Washnitzer--Monsky pour $X_0$ lisse affine, et $H_{\mathrm{DR}}(X/W)$ modulo la torsion quand $X_0$ a un relèvement propre et lisse $X$.
\end{conjecture*}
\end{encadre}
{\small Mathlib n'a pas de cohomologie $p$-adique des schémas. On se reportera aux notes de Grothendieck (\lot{7}{1}{16}).\par}

\begin{encadre}
\begin{theoreme*}[Grothendieck, pages 74 à 83, catégories polynomiales\piste{137-polynomial-categories-pd-tensor-categories}]\nonverifie
Les catégories polynomiales strictes, catégories à produits finis munies d'une sous-catégorie additive de flèches de degré un et d'une graduation des $\mathrm{Hom}$ par l'homogénéité sous les scalaires, sont équivalentes aux $\otimes$-pd-catégories, catégories monoïdales symétriques additives munies d'un foncteur gradué exponentiel $\Gamma^*$ et de flèches $\alpha_{p,q} : \Gamma^{pq} \to \Gamma^p\Gamma^q$ ; le passage inverse est $\mathbf P_p(X, Y) = \mathrm{Hom}(\Gamma^p X, Y)$.
\end{theoreme*}
\end{encadre}
{\small La page~83 affirme l'équivalence sans la démontrer : \q{Les deux notions sont donc équivalentes}. Mathlib n'a aucune des deux notions. L'équivalence complète, avec les produits, l'additivité, les multidegrés, le produit tensoriel strict et la formule exponentielle, n'est pas vérifiée en Lean ; son c\oe ur, la partie qui porte le passage inverse $\mathbf P_p(X, Y) = \mathrm{Hom}(\Gamma^p X, Y)$ et la règle de composition, l'est, et il est énoncé ci-dessous. On se reportera aux notes de Grothendieck (\lot{137}{4}{74}).\par}

\begin{encadre}
\begin{theoreme*}[D'après les pages 77 à 83, forme réduite par l'édition, le c\oe ur de l'équivalence\piste{137-polynomial-categories-pd-tensor-categories}]\verifie{Folder137}
Soit $\mathcal C$ une catégorie, celle des flèches de degré un. Appelons \emph{catégorie polynomiale sur $\mathcal C$} la donnée d'ensembles $\mathbf P_n(X, Y)$, d'une composition $\mathbf P_p(X, Y) \times \mathbf P_q(Y, Z) \to \mathbf P_{qp}(X, Z)$ associative, d'identités de degré un, et d'une bijection $\mathrm{Hom}_{\mathcal C}(X, Y) \simeq \mathbf P_1(X, Y)$ compatible à la composition ; elle est \emph{représentable} si chaque foncteur $Y \mapsto \mathbf P_n(X, Y)$ sur $\mathcal C$ est représenté par un objet $\Gamma^n X$ et une flèche universelle $\gamma_n \in \mathbf P_n(X, \Gamma^n X)$. Appelons \emph{structure à puissances divisées sur $\mathcal C$} la donnée de foncteurs $\Gamma^n : \mathcal C \to \mathcal C$, de transformations naturelles $\alpha_{p,q} : \Gamma^{pq} \to \Gamma^p\Gamma^q$, d'un isomorphisme $\varepsilon : \Gamma^1 \simeq \mathrm{id}$, satisfaisant le carré d'associativité de la page~80 et les conditions d'unité $(**)$. Alors :
\begin{enumerate}
\item d'une structure à puissances divisées, $\mathbf P_p(X, Y) = \mathrm{Hom}(\Gamma^p X, Y)$ avec $g \circ f = g \circ \Gamma^q(f) \circ \alpha_{q,p}$ fait une catégorie polynomiale représentable, par $\Gamma^n X$ et $\gamma_n = \mathrm{id}$ ; si $\mathcal C$ est additive, cette composition est additive en $g$ ;
\item d'une catégorie polynomiale représentable, les $\Gamma^n$ sont des foncteurs, les $\alpha_{p,q}$ définis par la factorisation de $\gamma_p \circ \gamma_q$ sont naturels, $\Gamma^1 \simeq \mathrm{id}$, et le carré de la page~80 et $(**)$ sont vrais ;
\item les deux passages sont inverses l'un de l'autre : le premier suivi du second redonne exactement $\Gamma$, $\alpha$ et $\varepsilon$ ; le second suivi du premier redonne la catégorie polynomiale de départ à isomorphisme près, par $\varphi \mapsto \varphi \circ \gamma_p$.
\end{enumerate}
\end{theoreme*}
\end{encadre}
{\small \emph{Démonstration.} Une structure à puissances divisées est une comonade graduée par le monoïde multiplicatif $(\mathbf N, \cdot)$, et la catégorie polynomiale qu'elle définit en est la catégorie de Kleisli. (1) L'associativité de la composition des pages~81 et~82 est ce qui demande la naturalité de $\alpha$ : dans $h \circ (g \circ f)$ et $(h \circ g) \circ f$, on fait passer $\Gamma^{rq}(f)$ à travers $\alpha_{r,q}$ par naturalité, et il reste à comparer $\alpha_{r,q} \circ \alpha_{rq,p}$ et $\Gamma^r(\alpha_{q,p}) \circ \alpha_{r,qp}$, qui est le carré de la page~80. Les unités sont $(**)$ et la naturalité de $\varepsilon$. La représentabilité est immédiate : $\varphi \circ \gamma_n = \alpha_{1,n} \circ \varepsilon \circ \varphi = \varphi$, par $(**)$. (2) Tout se lit à travers la bijection $\varphi \mapsto \varphi \circ \gamma_n$ de $\mathrm{Hom}(\Gamma^n X, Y)$ sur $\mathbf P_n(X, Y)$ : on définit $\Gamma^n(f)$ par $\Gamma^n(f) \circ \gamma_n = \gamma_n \circ f$, $\alpha_{p,q}$ par $\alpha_{p,q} \circ \gamma_{pq} = \gamma_p \circ \gamma_q$ (la page~80 : \q{ce qui revient au même}), et $\varepsilon$ par $\varepsilon \circ \gamma_1 = \mathrm{id}$. Chaque identité à démontrer, fonctorialité, naturalité, carré, unités, devient après composition avec le $\gamma$ convenable une égalité entre composés de $\gamma$, qui découle de l'associativité de $\mathbf P$ ; par exemple les deux côtés du carré deviennent $\gamma_p \circ \gamma_q \circ \gamma_r$, parenthésé de deux façons. L'inverse de $\varepsilon$ est la flèche de degré un qui correspond à $\gamma_1$. (3) Dans le premier passage, $\varphi \mapsto \varphi \circ \gamma_n$ est l'identité ; le second passage reconstruit donc $\Gamma$, $\alpha$ et $\varepsilon$ à partir de $\gamma_n = \mathrm{id}$, et les retrouve par $(**)$. Dans l'autre sens, la bijection $\varphi \mapsto \varphi \circ \gamma_p$ envoie $g \circ \Gamma^q(f) \circ \alpha_{q,p}$ sur $(g \circ \gamma_q) \circ (f \circ \gamma_p)$, par le même calcul. Rien n'utilise l'additivité, le produit tensoriel ni la formule exponentielle : le c\oe ur de l'équivalence vaut pour une catégorie $\mathcal C$ quelconque, et la règle de composition de la page~82 est forcée par la représentabilité. Les degrés sont portés par une équation $pq = n$, ce qui évite toute conversion d'indices. Lean (\href{\leanblob/Folder137.lean}{\texttt{Folder137.\allowbreak{}lean}}) définit les deux notions (\texttt{PDCategory}, \texttt{PolynomialCategory}, \texttt{PolynomialCategory.\allowbreak{}Representation}) et démontre (1) (\texttt{PDCategory.\allowbreak{}toPoly}, \texttt{PDCategory.\allowbreak{}representation}, \texttt{PDCategory.\allowbreak{}toPoly\_comp\_add}), (2) (\texttt{Representation.\allowbreak{}toPD}) et (3) (\texttt{PDCategory.\allowbreak{}toPD\_representation\_\(\Gamma\)\_map}, \texttt{PDCategory.\allowbreak{}toPD\_representation\_\(\alpha\)}, \texttt{PDCategory.\allowbreak{}toPD\_representation\_\(\varepsilon\)}, \texttt{Representation.\allowbreak{}toPDToPolyIso}) ; un modèle trivial, tous les $\Gamma^n$ égaux à l'identité, montre que les axiomes sont compatibles (\texttt{PDCategory.\allowbreak{}trivial}). Ne sont pas démontrés : les axiomes a) à c) (produits, additivité de $\mathbf P_1$), la décomposition d) par multidegrés et l'homogénéité sous les scalaires, f) et g), la moitié $\mathbf 1 \simeq \Gamma^q(\mathbf 1)$ de $(***)$, le produit tensoriel strict, la formule exponentielle $\Gamma^*(X \oplus Y) \simeq \Gamma^*X \otimes \Gamma^*Y$ et l'exemple des $k$-modules (théorème de Roby).\par}
